%%%%%%%%%%  scan idx 217  |  folio 204  %%%%%%%%%%
\SGAPrintPage{204}%
\[
  i_s^*(P^n_{X/S})\simeq P^n_{X_s},
\]
whence a restriction homomorphism
\[
  i_s^*:\Diff^n(E,F)
  \longrightarrow
  \Diff^n\bigl(i_s^*(E),i_s^*(F)\bigr),
\]
for locally free $\underline{O}_X$-Modules $E,F$ of finite type (and an analogous homomorphism for operators with complex coefficients). Let $E$ be a complex, with bounded degrees, of differential operators with complex coefficients on $X$; for $E$ to be elliptic (1.2.5), it is necessary and sufficient that, for every $s\in S$, $i_s^*(E)$ be an elliptic complex (verification left to the reader).

\sgaStableTarget{sga6.e02.sec.2.occ2}\subsection*{2. Rigidity lemmas}

\sgaStableTarget{sga6.e02.para.2-0}2.0. It is \textquotedblleft well known\textquotedblright{} that compact differentiable varieties are \textquotedblleft rigid\textquotedblright, i.e. admit no continuous deformation other than a trivial one. However, since this result will be essential to us in \S3, and references seem rather scarce, it may be useful to give precise statements and a few proofs.

We fix a topological space $S$. The mixed varieties considered will be assumed to be of class $C^r$, with $r\ge1$ (possibly $r=\infty$).

\sgaStableTarget{sga6.e02.para.2-1}\textbf{2.1. Notation.} Let $X,Y$ be mixed $S$-varieties. We denote by $\Homop(X,Y)$ the set of $S$-morphisms from $X$ to $Y$, and by $\underline{\Homop}(X,Y)$ the sheaf on $X$ of germs of $S$-morphisms from $X$ to $Y$ (defined by
\[
  \underline{\Homop}(X,Y)(U)=\Homop(U,Y)
\]
for $U$ open in $X$).

Let $s\in S$. We denote by $i_s$ the inclusion of $X_s$ into $X$. There is a canonical induction homomorphism
\sgaStableTarget{sga6.e02.eq.2-1-1}\begin{equation}\tag{2.1.1}
  i_s^{-1}\underline{\Homop}(X,Y)
  \longrightarrow
  \underline{\Homop}(X_s,Y_s),
\end{equation}
%%%%%%%%%%  scan idx 218  |  folio 205  %%%%%%%%%%
\SGAPrintPage{205}%
which is an epimorphism, as we see immediately. When $Y=S\times\underline{\underline{R}}$, the sheaf $\underline{\Homop}(X,Y)$ is canonically identified with $\underline{O}_X$, and the epimorphism (2.1.1) with the epimorphism
\[
  i_s^{-1}(\underline{O}_X)\longrightarrow i_s^*(\underline{O}_X)=\underline{O}_{X_s}
\]
considered in (1.1).

\sgaStableTarget{sga6.e02.prop.2-2}\textbf{Proposition}\SGAAuthorFootnote{1}{Communicated to the editor by A. Douady.} \textbf{2.2.} \uline{Let} $X,Y$ \uline{be mixed} $S$\uline{-varieties, and} $s\in S$\uline{. Suppose that the fiber} $X_s$ \uline{is paracompact and that} $s$ \uline{has a fundamental system of paracompact neighborhoods. Let} $A$ \uline{be a closed subset of} $X_s$\uline{, and} $f\in H^0(A,\underline{\Homop}(X_s,Y_s))$\uline{.
The subsheaf} $F$ \uline{of} $i_s^{-1}\underline{\Homop}(X,Y)$ \uline{formed by the germs \textquotedblleft that extend} $f$\uline{\textquotedblright{} (i.e. whose image under} (2.1.1) \uline{coincides with} $f$ \uline{over} $A$\uline{) is soft.}

\uline{Proof.} Let $X'$ be an open neighborhood of $x$ in $X$, isomorphic (as a mixed variety) to $S'\times U$, where $S'$ is a paracompact neighborhood of $s$ and $U$ a relatively compact open subset of $\underline{\underline{R}}^m$, and such that the image under $f$ of $X'_s\cap A$ is contained in an open subset of $Y$ isomorphic to $S'\times V$, where $V$ is a convex open subset of $\underline{\underline{R}}^n$ containing the origin. By (T.F. II 3.4.1), it suffices to show that every section of $F$ over a closed subset $A_0$ of $X_s$ contained in $X'_s$ extends to $X'_s$. Now $X'_s$ ($\simeq\{s\}\times U$) is relatively compact, so a closed subset of $X_s$ contained in $X'_s$ is compact. Identify $X'$ (resp. $Y'$) with $S'\times U$ (resp. $S'\times V$), put $A'=A\cap X'_s$, and let $A_0$ be a compact subset of $X'_s=\{s\}\times U$. Since $A'$ is closed in $\{s\}\times U$, which is paracompact, we have ([T.F.]~II 3.3.1)\SGAAuthorFootnote{1}{A closed subset of a paracompact space has a fundamental system of paracompact neighborhoods.}
\[
  H^0\bigl(A',\underline{\Homop}(X_s,Y_s)\bigr)
  =\varinjlim_{U'\supset A'}\Homop(U',Y'_s),
\]
over the open subsets $U'$ of $X'_s$ containing $A'$. Similarly, since $A_0$ is closed in $S'\times \overline U$, which is paracompact (a product of a paracompact space and a compact space), we have
\[
  H^0\bigl(A_0,i_s^{-1}\underline{\Homop}(X,Y)\bigr)
  =\mathop{\underrightarrow{\mathrm{Lim}}}\limits_{U_0\supset A_0}\Homop(U_0,Y'),
\]
over the open subsets $U_0$ of $X'$ containing $A_0$. The proposition will therefore be a
%%%%%%%%%%  scan idx 219  |  folio 206  %%%%%%%%%%
\SGAPrintPage{206}%
consequence of the following lemma:

\sgaStableTarget{sga6.e02.lem.2-2-1}\textbf{Lemma 2.2.1.} \uline{Let} $U$ \uline{be an open subset of} $\underline{\underline{R}}^m$, $V$ \uline{a convex open subset of} $\underline{\underline{R}}^n$ \uline{containing} $0$\uline{. Let} $A_0,A_1$ \uline{be closed subsets of} $\{s\}\times U$\uline{, let} $U_0$ (resp. $U_1$\uline{) be an open subset of} $S\times U$ (resp. $\{s\}\times U$\uline{) containing} $A_0$ (resp. $A_1$\uline{). Let} $f_1:U_1\to V$ \uline{be a morphism and let} $f_0:U_0\to S\times V$ \uline{be an} $S$\uline{-morphism inducing on} $U_0\cap(\{s\}\times U)$ \uline{a morphism defining the same germ as} $f_1$ \uline{in a neighborhood of} $A_1$\uline{. There then exists an open subset} $U'$ \uline{of} $S\times U$ \uline{containing} $s\times U$\uline{, and an} $S$\uline{-morphism} $\bar f:U'\longrightarrow S\times V$ \uline{such that}
\begin{itemize}
\item[(i)] \uline{the germ of} $\bar f$ \uline{along} $\{s\}\times U$ \uline{extends the germ defined by} $f_0$;
\item[(ii)] \uline{the morphism induced by} $\bar f$ \uline{on} $\{s\}\times U$ \uline{defines the same germ as} $f_1$ \uline{in a neighborhood of} $A_1$.
\end{itemize}

\uline{Proof.} Replacing $f_1$ by $af_1$, where $a$ is a $C^r$ function on $U$, equal to $1$ in a neighborhood of $A_1$ and with support in $U_1$, we may assume that $U_1=U$. Let $b$ be a $C^r$ function on $U$, equal to $1$ on $A_0$ and with support in $U_0\cap(\{s\}\times U)$. Then
\[
  U'=U_0\cup(S\times(U-\operatorname{Supp}b))
\]
is an open subset of $S\times U$ containing $\{s\}\times U$. The morphism $\bar f:U'\to S\times V$ defined by
\[
  \bar f(t,x)=\bigl(t,\ b(x)\pr_2f_0(t,x)+(1-b(x))f_1(x)\bigr)
\]
visibly satisfies conditions (i) and (ii).

% Author-claim audit ASTRA-II-AUTHOR-CLAIM-206 is recorded separately.
% The source statement and proof are retained in both English forms.
\sgaStableTarget{sga6.e02.cor.2-3}\textbf{Corollary 2.3.} \uline{Under the hypotheses of} (2.2)\uline{, the sheaves} $i_s^{-1}\underline{\Homop}(X,Y)$, $\underline{\Homop}(X_s,Y_s)$ \uline{are soft. In particular, the sheaves} $i_s^{-1}(\underline{O}_X)$, $\underline{O}_{X_s}$ \uline{are soft}\SGAAuthorFootnote{1}{For $\underline{O}_{X_s}$, this is \textquotedblleft very well known\textquotedblright. The editor does not know whether $\underline{O}_X$ is soft when the hypotheses of (2.2) hold for every $s\in S$ and $X$ is paracompact. Under those hypotheses, this would trivially imply the softness of $i_s^{-1}(\underline{O}_X)$ as well as (2.5) below.}\uline{. For every closed subset} $A$ \uline{of} $X_s$\uline{, the map}
\[
  H^0\bigl(X_s,i_s^{-1}\underline{\Homop}(X,Y)\bigr)
  \longrightarrow
  H^0\bigl(A,\underline{\Homop}(X_s,Y_s)\bigr)
\]
\uline{deduced from} (2.1.1) \uline{is surjective.}
%%%%%%%%%%  scan idx 220  |  folio 207  %%%%%%%%%%
\SGAPrintPage{207}%
\uline{Proof.} The softness of $i_s^{-1}\underline{\Homop}(X,Y)$ follows from (2.2) applied with $A=\varnothing$; taking $S=s$, we deduce that $\underline{\Homop}(X_s,Y_s)$ is soft. The last assertion follows from (2.2).

\sgaStableTarget{sga6.e02.cor.2-4}\textbf{Corollary 2.4.} \uline{Under the hypotheses of} (2.2)\uline{, let} $E$ \uline{be an} $\underline{O}_X$\uline{-Module. For every closed subset} $A$ \uline{of} $X_s$\uline{, the canonical homomorphism}
\[
  H^0\bigl(X_s,i_s^{-1}(E)\bigr)
  \longrightarrow
  H^0\bigl(A,i_s^*(E)\bigr),
\]
\uline{deduced from the epimorphism}
\[
  i_s^{-1}(E)\longrightarrow i_s^*(E)
  =i_s^{-1}(E)\otimes_{i_s^{-1}(\underline{O}_X)}\underline{O}_{X_s},
\]
\uline{is an epimorphism.}

\uline{Proof.} Since $i_s^{-1}(\underline{O}_X)$ is soft (2.3), the functor $H^0(A,.)$ is exact on the category of $i_s^{-1}(\underline{O}_X)$-Modules (T.F. II 3.7.1, 3.4.2, 4.4.3), so the homomorphism
\[
  H^0\bigl(A,i_s^{-1}(E)\bigr)
  \longrightarrow
  H^0\bigl(A,i_s^*(E)\bigr)
\]
is an epimorphism. Since $i_s^{-1}(E)$ is soft, the homomorphism
\[
  H^0\bigl(X_s,i_s^{-1}(E)\bigr)
  \longrightarrow
  H^0\bigl(A,i_s^{-1}(E)\bigr)
\]
is an epimorphism, and we are done.

\sgaStableTarget{sga6.e02.cor.2-5}\textbf{Corollary 2.5.} \uline{Let} $p:X\to S$ \uline{be a mixed variety, and} $s\in S$\uline{. Suppose that} $s$ \uline{has a fundamental system of paracompact neighborhoods} $S'$ \uline{such that the} $p^{-1}(S')$ \uline{form a fundamental system of paracompact neighborhoods of} $X_s$ \uline{(a condition satisfied, for example, when} $S$ \uline{is locally compact and} $p$ \uline{is a proper map of topological spaces). Then, for every} $\underline{O}_X$\uline{-Module} $E$\uline{, we have}
\[
  R^i p_*(E)_s=0 \quad \text{for } i>0.
\]
%%%%%%%%%%  scan idx 221  |  folio 208  %%%%%%%%%%
\SGAPrintPage{208}%
\uline{Proof.} By (T.F. II 4.17.1), we have
\[
  R^i p_*(E)_s=H^i\bigl(X_s,i_s^{-1}(E)\bigr).
\]
But, since $i_s^{-1}(\underline{O}_X)$ is soft (2.3), we have
\[
  H^i\bigl(X_s,i_s^{-1}(E)\bigr)=0\quad \text{for } i>0,
\]
whence the assertion.

\sgaStableTarget{sga6.e02.cor.2-6}\textbf{Corollary 2.6.} \uline{Let} $X,Y$ \uline{be mixed} $S$\uline{-varieties, and} $s\in S$\uline{. Suppose that} $X$ \uline{satisfies the hypotheses of} (2.5)\uline{. For every morphism} $f:X_s\to Y_s$\uline{, there exists an open neighborhood} $S'$ \uline{of} $s$ \uline{and a morphism} $\bar f:X|S'\longrightarrow Y|S'$ \uline{such that} $\bar f_s=f$.

\uline{Proof.} This follows immediately from (2.3) (with $A=X_s$) and the fact that, by the hypotheses on $X$, we have
\[
  H^0\bigl(X_s,i_s^{-1}\underline{\Homop}(X,Y)\bigr)
  =\operatorname*{Lim}_{S'\ni s}\Homop(X|S',Y)\qquad (S'\ \text{open})
\]
(T.F. II 3.3.1).

\sgaStableTarget{sga6.e02.prop.2-7}\textbf{Proposition 2.7.} \uline{Let} $X,Y$ \uline{be mixed} $S$\uline{-varieties,} $f:X\to Y$ \uline{a morphism,} $s\in S$, $x\in X_s$\uline{. If the tangent linear map of} $f_s:X_s\to Y_s$ \uline{is invertible at} $x$\uline{, then} $f$ \uline{is a local isomorphism at} $x$\uline{, i.e. there exists an open neighborhood} $U$ \uline{of} $x$ \uline{in} $X$ \uline{and an open subset} $V$ \uline{of} $Y$ \uline{such that} $f$ \uline{induces an isomorphism from} $U$ \uline{onto} $V$.

\uline{Proof.} This is an easy consequence of the implicit function theorem. The details are left to the reader.

\sgaStableTarget{sga6.e02.cor.2-8}\textbf{Corollary 2.8.} \uline{Let} $X,Y$ \uline{be mixed} $S$\uline{-varieties, and} $s\in S$\uline{. Suppose that} $s$ \uline{has a fundamental system of neighborhoods} $S'$ \uline{such that the} $X|S'$ \uline{(resp.} $Y|S'$\uline{) form a fundamental system of neighborhoods of} $X_s$ \uline{(resp.} $Y_s$\uline{) (a condition satisfied, for example, when} $S$ \uline{is locally compact and} $X$ \uline{(resp.} $Y$\uline{) is proper over} $S$\uline{). Let} $f:X\to Y$ \uline{be an} $S$\uline{-morphism such that} $f_s:X_s\to Y_s$ \uline{is an isomorphism. There then exists a neighborhood} $S'$

